\subsection{Automatic nonzero norms by polynomial-adic division}
\label{sec:adic-squarefree-compact}

Let \(K\) be an exact field with a primitive-recursive arithmetic
presentation.  Let \(f\in K[X]\) have degree \(d\geq1\), and suppose
\(K[X]/(f)\) is reduced.  A sufficient, effectively checkable input is
a Bezout certificate \(Uf+Vf'=1\); over \(\mathbb Q\) this is the
ordinary squarefree condition.  Over imperfect fields, the displayed
certificate means separability, which must not be confused with the
weaker ring-theoretic absence of repeated irreducible factors.

\begin{theorem}[The polynomial-adic norm certificate]
\label{thm:adic-squarefree-compact}
Let \(m\geq1\), and let
\(G=\sum_{j<m}A_j(X)Y^j\ne0\) be reduced in \(Y\).
Then \(N=N_{Y^m-f}(G)\ne0\), and
\[
 \deg N\leq W:=\max_{A_j\ne0}(m\deg A_j+jd).
\]
More precisely, bounded division computes an index \(k<m\), an integer
\(a\geq0\), and a polynomial \(B\not\equiv0\pmod f\), with
\(s=ma+k\), such that the exact quotient \(H=N/f^s\) exists and
\begin{equation}
 \label{eq:adic-exact-residue}
 H\equiv(-1)^{k(m-1)}B^m\pmod f,
 \qquad H\bmod f\ne0.
\end{equation}
Neither \(\gcd(m,d)=1\) nor a rational root of \(f\) is assumed.
\end{theorem}

\begin{proof}
For each nonzero \(A_j\), repeatedly divide by \(f\) while the
remainder is zero.  Since a successful division reduces the degree by
\(d\geq1\), at most \(\deg A_j+1\) divisions give
\[
 A_j=f^{a_j}B_j,\qquad B_j\bmod f\ne0,
 \qquad a_jd\leq\deg A_j.
\]
The weights \(w_j=ma_j+j\) have distinct residues modulo \(m\), so
their minimum is attained at a unique index \(k\).  Set
\(a=a_k\), \(B=B_k\), and \(s=w_k\).

Expand the multiplication determinant in the basis
\(1,Y,\ldots,Y^{m-1}\).  In column \(i\), the coefficient
\(A_jY^j\) contributes to row \(r\equiv i+j\pmod m\), with entry
\(A_j f^q\), where \(i+j=r+mq\) and \(q\in\{0,1\}\).
For a determinant permutation \(r=\sigma(i)\), the selected indices
\(j_i\) and carries \(q_i\) therefore satisfy
\[
 \sum_i j_i=m\sum_iq_i,\qquad
 \sum_i a_{j_i}+\sum_iq_i=\frac1m\sum_iw_{j_i}.
\]
The latter integer is at least \(s\), and exceeds \(s\) if any
\(j_i\ne k\).  Terms selecting a zero coefficient vanish.
The unique term selecting \(k\) in every column is the cyclic
permutation \(i\mapsto i+k\pmod m\), of sign
\((-1)^{k(m-1)}\), and its total carry is \(k\).
Thus, by factoring the displayed powers from the finite determinant sum,
one obtains an actual polynomial identity
\[
 N=f^s\bigl((-1)^{k(m-1)}B^m+fQ\bigr).
\]
This constructs the exact division quotient
\(H=(-1)^{k(m-1)}B^m+fQ\), and proves its residue formula.
In the reduced ring \(K[X]/(f)\), the nonzero residue of \(B\) has
nonzero \(m\)-th power.  Hence \(H\bmod f\ne0\), so \(H\ne0\).
Because \(K[X]\) is a domain and \(f\ne0\), also \(N\ne0\).

For the degree bound, every nonzero determinant term has degree at most
\[
 \sum_i\deg A_{j_i}+d\sum_iq_i
 =\frac1m\sum_i(m\deg A_{j_i}+dj_i)\leq W.
\]
Addition cannot increase this bound.  Furthermore
\(sd=(ma+k)d\leq m\deg A_k+kd\leq W\), so the quotient has
degree at most \(W-sd\).  If \(D=\max_j\deg A_j\), the selected
exponent has the explicit bound \(s\leq mD+m-1\).
The integer-count equalities involving \(1/m\) do not divide by an
element of \(K\), so characteristics dividing \(m\) cause no problem.
\end{proof}

\paragraph{The reducedness certificate, without factorization.}
Suppose \(Uf+Vf'=1\), and \(f\mid b^r\), with \(r\geq1\).
Compute \(g=\gcd(f,b)\), \(f=gh\), \(b=gq\), and
\(Af+Bb=g\).  Cancelling \(g\ne0\) gives \(Ah+Bq=1\).
Differentiating \(f=gh\) in the squarefree certificate gives
\(Ch+Dg=1\), where \(C=Ug+Vg'\), \(D=Vh'\).
Multiplying these two identities gives
\[
 1=Eh+Fb,
 \quad E=CAh+CBq+DAg,
 \quad F=DB.
\]
Write \(b^r=hR\).  The finite geometric identity now gives
\[
 1=(Fb)^r+(1-Fb)\sum_{i<r}(Fb)^i
   =h\left(F^rR+E\sum_{i<r}(Fb)^i\right)=hJ.
\]
Therefore \(b=gq=gh(qJ)=f(qJ)\).  This is an explicit witness that
\(f\mid b\), proving reducedness.  After this one lemma, the norm
producer uses no derivative inverse, formal root, or extension-field
construction.

\paragraph{Concrete gain.}
For \(m=2\), \(f=X^4+1\), and \(G=Y-X^2\) over \(\mathbb Q\),
the global degree weights tie at \(4\), \(\gcd(m,d)=2\), and
\(f\) has no rational root.  The polynomial-adic weights are \(0,1\),
and \eqref{eq:adic-exact-residue} detects
\(N=X^4-(X^4+1)=-1\).
The residue algebra need not be a field: for \(f=X^2-1\),
\(G=(X-1)+Y\), the least residue \(u-1\) is a nonzero zero divisor,
but its square is \(2-2u\ne0\), and the norm is \(2-2X\).

\paragraph{iPRA and Lean obligations.}
The implementation consists of bounded coefficient scans, Euclidean
division, a degree-bounded gcd computation with Bezout witnesses,
at most \(\deg A_j+1\) repeated divisions per row, and a
division-free determinant computation bounded by \(m!\).
The proof records the exact quotient identity, the carry-sum identity,
and the finite geometric identity as quantifier-free loop invariants.
Nonzero residue yields nonzero norm, and the degree cap \(W\) then
bounds a scan returning an actual nonzero coefficient.  Normalize a
divisor for remainder computation while retaining its leading-coefficient
unit in the proved identities.  Never invert an arbitrary nonzero residue
of \(K[X]/(f)\).  The corresponding Lean endpoint must export an
object-iPRA derivation for these fixed primitive-recursive programs;
an abstract-field Lean theorem alone is a separate semantic specification.
